Doleo etiam eos non adesse, quos maxime oportuit\footnote{M oportet.}, Qui cum et privatim et publice toties me crimine haereseos\footnote{M in haereseos.} prophanarunt, nunc cum instet cognitio causae se subtraxerunt, haereticae pravitatis inquisitores dico, qui fraternam monitionem\footnote{M Lücke, Mc paternam admonitionem.} et doctrinam postposuerunt criminationibus suis.

[255] Dominus doctor Eccius\footnote{U Scopus. De potestate, immo de primatu Romani pontificis (Uv a. R. Additio). Eccii protestatio, Mc [rot] = U.}.

In nomine tuo, dulcis Jesu\footnote{fehlt in M.}.

Antequam in harenam descendam protestor coram, vobis illustrissimis, nobilibus, magnificis et excellentissimis dominis omnia per me dicenda et dicta, inprimis esse subiecta iudicio primae sedis et domini in ea sedentis, deinde\footnote{M dehinc.} quorumlibet aliorum, quorum est errantes corrigere et ad veritatis cognitionem reducere. Et quia reverendus pater in praefatione sua, quasi se excusando, testatur se libenter 'hanc' obmisisse\footnote{U omisisse.} materiam pro summi pontificis reverentia, nisi per meam pertractus fuisset propositionem: at meminerit reverendus pater: nisi prius ipse in resolutorio ante tempora Sylvestri negasset Romanum\footnote{M Romanam.} pontificem aliis\footnote{M aliis esse.} superiorem, non fuisset necessarium me tredecimam illam posuisse propositionem. Et quod in actis coram legato\footnote{M legatis.} sedis apostolicae beatum Pelagium pontificem evangelica scripta\footnote{M dicta.} torsisse causatur. Qui tamen maxime omnium, ad sanctorum patrum intentionem verba Christi accepit. Frustra ergo reverendus pater in me causam torquere nititur, cui ipse non semel occasionem praestitit: sed iam ambagibus resectis\footnote{P reiectis.} principale institutum, deo duce, aggrediamur. Reverende pater, vestra conclusio XIII. meae contradicens, fatetur Romanam ecclesiam' esse' aliis superiorem iuxta frigidissima Romanorum pontificum decreta, intra quadringentos annos\footnote{M quadringenta annis.} nata. Quibus obstare dicitis, textum sacrae scripturae et historias approbatas mille et\footnote{et fehlt in M.} centum 'annorum'. Contra quam sic oppono.

Monarchia et unus principatus in ecclesia dei est de iure divino et a Christo institutus. Quare textus sacrae scripturae vel historiae approbatae ei non adversantur, quoniam ecclesia illa militans, quae est velut unum corpus iuxta divi Pauli sententiam, est instituta et facta ad imaginem ecclesiae triumphantis: in qua est una monarchia
% empalme: la página sigue "omni-" + "bus" = "omnibus" (ver test-p02) -->
